\newpage
\appendix
\onecolumn

\icmltitle{SuiteMerge: Deconstructing the Task Suite Bias in Model Merging\\--- Supplementary Material ---}

\section{Additional Mathematical Derivations and Model II Landscape Specifications}
\label{app:model_ii}

In this section, we provide the full mathematical specification of our calibrated non-convex, coupled Model~II sensitivity landscape used in the simulation experiments.

\subsection{Coupled Model II Landscape Equations}
For a task suite $S$ and merging coefficient parameters $\theta_S$, let $\alpha_k(l)$ denote the merging coefficient for task $k \in S$ at layer $l \in \{1, \dots, L\}$.
The local sensitivity loss $\mathcal{S}_k^{(l)}(\alpha_k(l))$ is governed by:
\begin{equation}
\mathcal{S}_k^{(l)}(\alpha_k(l)) = A_k^{(l)} \left( \alpha_k(l) - \alpha_{k, \text{opt}}^{(l)} \right)^2 + B_k^{(l)} \left( \alpha_k(l) - \alpha_{k, \text{opt}}^{(l)} \right)^4 + \text{Interference}_k^{(l)}(\theta_S)
\end{equation}
where the coupling/interference term is mathematically modeled as:
\begin{equation}
\text{Interference}_k^{(l)}(\theta_S) = \sum_{k' \in S, k' \neq k} D_{k, k'}^{(l)} \left( \alpha_k(l) - \alpha_{k'}(l) \right)^2
\end{equation}

The coefficients $A_k^{(l)}$ and $B_k^{(l)}$ govern the local quadratic and quartic curvature, respectively, which correspond to the sensitivity of layer $l$ to modifications of its weights.
The term $D_{k, k'}^{(l)}$ is the pairwise representational conflict penalty at layer $l$ between task $k$ and task $k'$.

\subsection{Parameter Calibration and Empirical Protocol}
To ensure that the simulator represents a physically plausible model-merging landscape, the curvature parameters, optimal coefficients, and representational conflict penalties were calibrated against empirical Vision Transformer (ViT-B/32) classification statistics.

\paragraph{Empirical Calibration Protocol:} 
The exact empirical protocol used to extract the classification statistics for calibrating the quadratic and quartic sensitivity parameters $A_k^{(l)}$ and $B_k^{(l)}$ is structured as follows:
\begin{enumerate}
    \item \textbf{Expert Fine-tuning:} Starting from a pre-trained ViT-B/32 backbone (86M parameters), task-specific experts were fine-tuned independently on MNIST, FashionMNIST, CIFAR-10, and SVHN, respectively.
    \item \textbf{Layer-Wise Coefficient Perturbation:} At each individual layer $l \in \{1, \dots, 12\}$, we systematically perturbed the task-specific merging coefficient $\alpha_k(l)$ away from the default uniform merging point ($0.5$) in incremental steps of $0.05$ across the physically valid interval $[0.0, 1.0]$.
During this layer-wise sweep, all other layers $l' \neq l$ were held static at the uniform baseline coefficient of $0.5$.
    \item \textbf{Validation Set Evaluation:} For each step of the perturbation sweep at layer $l$, we evaluated the classification accuracy of the resulting model on the validation set of task $k$.
This yielded an empirical layer-wise accuracy-sensitivity curve $f_k^{(l)}(\alpha_k(l))$ for each layer and task.
    \item \textbf{Mathematical Fit and Parameter Extraction:} To extract the quadratic and quartic sensitivity parameters, we fit a fourth-degree polynomial of the form $g(\alpha_k(l)) = - A_k^{(l)} \left( \alpha_k(l) - \alpha_{k, \text{opt}}^{(l)} \right)^2 - B_k^{(l)} \left( \alpha_k(l) - \alpha_{k, \text{opt}}^{(l)} \right)^4 + C_k^{(l)}$ to the empirical accuracy data $f_k^{(l)}(\alpha_k(l))$ via least-squares regression, where $\alpha_{k, \text{opt}}^{(l)}$ is the empirical optimal merging coefficient (peak accuracy) and $C_k^{(l)}$ is a scaling constant. 
    \item \textbf{Observations and Simulation Mapping:} Our regression analyses yielded two primary observations:
    \begin{itemize}
        \item Early layers ($l \le 3$, responsible for basic feature extraction and alignment) and late layers ($l \ge 10$, responsible for task-specific classification head alignment) exhibited steep, sharp sensitivity landscapes, indicating high sensitivity to merging coefficient perturbation.
These layers are mapped in our simulator by sampling $A_k^{(l)} \sim \text{Uniform}(0.8, 1.2)$.
        \item Middle layers ($l \in \{4, \dots, 9\}$, which capture more general and transferrable representations) were highly robust to perturbation, exhibiting broad, shallow sensitivity landscapes.
These layers are mapped in our simulator by sampling $A_k^{(l)} \sim \text{Uniform}(0.2, 0.4)$.
        \item Across all layers, we observed a rapid performance drop-off when coefficients drifted significantly beyond the optimal basin, which is modeled in our quartic sensitivity coefficient by setting $B_k^{(l)} = 0.5 \cdot A_k^{(l)}$.
    \end{itemize}
\end{enumerate}

Based on these empirical calibration findings, the remaining simulator parameters are configured as:
\begin{itemize}
    \item \textbf{Ground-Truth Optimal Coefficients ($\alpha_{k, \text{opt}}^{(l)}$):} To simulate smooth trajectories across layers, we model $\alpha_{k, \text{opt}}^{(l)}$ as:
    \begin{itemize}
        \item Task 1 (MNIST): Linear increasing, $\alpha_{1, \text{opt}}^{(l)} = 0.3 + 0.4 \cdot \left(\frac{l-1}{L-1}\right)$.
        \item Task 2 (FashionMNIST): Linear decreasing, $\alpha_{2, \text{opt}}^{(l)} = 0.7 - 0.4 \cdot \left(\frac{l-1}{L-1}\right)$.
        \item Task 3 (CIFAR-10): Quadratic concave, $\alpha_{3, \text{opt}}^{(l)} = 0.2 + 1.2 \cdot \left(\frac{l-1}{L-1}\right) - 1.2 \cdot \left(\frac{l-1}{L-1}\right)^2$.
        \item Task 4 (SVHN): Constant trajectory, $\alpha_{4, \text{opt}}^{(l)} = 0.55$.
    \end{itemize}
    Minor high-frequency perturbations are added: $\alpha_{k, \text{opt}}^{(l)} \leftarrow \text{Clip}(\alpha_{k, \text{opt}}^{(l)} + \epsilon^{(l)}, 0.0, 1.0)$ with $\epsilon^{(l)} \sim \mathcal{N}(0, 0.02)$ to reflect local architectural fluctuations in actual neural networks.
    \item \textbf{Interference Matrix ($D_{k, k'}^{(l)}$):} The conflict scale is layer-dependent and is determined by the task suite's empirical domain relationships:
    \begin{itemize}
        \item \textbf{Suite A (Homogeneous):} $D_{1, 2}^{(l)} \sim \text{Uniform}(0.01, 0.03)$, reflecting high representational compatibility between simple grayscale shapes.
        \item \textbf{Suite B (Heterogeneous):} $D_{3, 4}^{(l)} \sim \text{Uniform}(0.20, 0.30)$, representing the severe representational clashing between RGB natural objects (CIFAR-10) and street numbers (SVHN) sharing the same visual features.
        \item \textbf{Suite C (Cross-Domain Digits):} $D_{1, 4}^{(l)} \sim \text{Uniform}(0.12, 0.18)$.
        \item \textbf{Suite D (Cross-Domain Objects):} $D_{2, 3}^{(l)} \sim \text{Uniform}(0.15, 0.21)$.
    \end{itemize}
\end{itemize}

\subsection{Surrogate Test-Time Adaptation Landscape}
To simulate the unsupervised prediction entropy minimization objective of online TTA under local batch bias, we add a non-convex cosine penalty to the sensitivity loss:
\begin{equation}
\mathcal{L}_{\text{TTA}}(\theta_S) = \sum_{k \in S} \left[ \sum_{l=1}^L \mathcal{S}_k^{(l)}(\alpha_k(l)) + \lambda_{\text{rug}} \cos(F \cdot \alpha_k(l)) \right]
\end{equation}
where $\lambda_{\text{rug}} = 0.05$ represents the ruggedness strength (non-convexity) and $F = 15.0$ represents the high-frequency oscillation factor of the prediction entropy surface.

---

\section{Detailed Physical Weight-Space Experimental Setup}
\label{app:physical_setup}

To confirm the empirical validity of our findings, we trained and merged physical neural network checkpoints.
Here, we outline the exact experimental and training hyperparameters.

\subsection{DeepCNN Model Architecture}
The physical weight-space validation is conducted on a custom 5-layer Convolutional Neural Network (\texttt{DeepCNN}) implemented in PyTorch:
\begin{enumerate}
    \item \textbf{Layer 1:} \texttt{Conv2d(in\_channels=1 or 3, out\_channels=16, kernel\_size=3, padding=1)} $\rightarrow$ \texttt{BatchNorm2d(16)} $\rightarrow$ \texttt{ReLU} $\rightarrow$ \texttt{MaxPool2d(kernel\_size=2)}
    \item \textbf{Layer 2:} \texttt{Conv2d(16, 32, kernel\_size=3, padding=1)} $\rightarrow$ \texttt{BatchNorm2d(32)} $\rightarrow$ \texttt{ReLU} $\rightarrow$ \texttt{MaxPool2d(kernel\_size=2)}
    \item \textbf{Layer 3:} \texttt{Conv2d(32, 64, kernel\_size=3, padding=1)} $\rightarrow$ \texttt{BatchNorm2d(64)} $\rightarrow$ \texttt{ReLU} $\rightarrow$ \texttt{MaxPool2d(kernel\_size=2)}
    \item \textbf{Layer 4 (Fully Connected):} \texttt{Linear(576, 128)} $\rightarrow$ \texttt{ReLU} $\rightarrow$ \texttt{Dropout(p=0.1)}
    \item \textbf{Layer 5 (Task Heads):} Separate linear heads are attached to the output of Layer 4 for each task:
    \begin{itemize}
        \item MNIST classification head: \texttt{Linear(128, 10)}
        \item FashionMNIST classification head: \texttt{Linear(128, 10)}
    \end{itemize}
\end{enumerate}
The model is modified to support layer-wise scalar merging coefficients.
Specifically, for any layer weight $W^{(l)}_{\text{merged}}$:
\begin{equation}
W^{(l)}_{\text{merged}} = \alpha^{(l)}_1 W^{(l)}_1 + \alpha^{(l)}_2 W^{(l)}_2
\end{equation}
where $W^{(l)}_1$ and $W^{(l)}_2$ are the corresponding weights of the MNIST and FashionMNIST expert networks, respectively.

\subsection{Expert Network Training Details}
We trained experts under two initialization and training regimes to evaluate weight-space connectivity:
\begin{enumerate}
    \item \textbf{Regime A: Disjoint Loss Basins (Scratch-Trained, Independent Initialization):}
    The MNIST expert and FashionMNIST expert are trained independently starting from separate random weight initializations (MNIST expert seeded with $42$, FashionMNIST expert seeded with $100$).
This ensures the experts converge to disjoint basins in the parameter space.
    \item \textbf{Regime B: Connected Loss Basins (Pre-trained Shared Initialization):}
    A base model is first pre-trained on a joint subset of MNIST and FashionMNIST for 2 epochs.
The MNIST and FashionMNIST experts are then fine-tuned from this shared pre-trained checkpoint on their respective datasets, placing them within a shared, linearly connected loss basin.
\end{enumerate}

All experts are trained using the Adam optimizer with a learning rate of $1 \times 10^{-3}$, a weight decay of $1 \times 10^{-4}$, and a batch size of 64 for 5 epochs.
Under these parameters, the individual expert models achieve high task-specific classification performance:
\begin{itemize}
    \item \textbf{MNIST Expert:} $98.40\%$ accuracy on MNIST test.
    \item \textbf{FashionMNIST Expert:} $88.10\%$ accuracy on FashionMNIST test.
\end{itemize}

---

\section{Nelder-Mead Solver Configuration}
\label{app:solver_config}

For both the simulated and physical experiments, our proposed \textbf{OFS-Tune} uses the Nelder-Mead simplex algorithm to optimize the continuous polynomial trajectory parameters $\theta_S$.
Here we detail the optimizer's configurations:

\begin{itemize}
    \item \textbf{Dimensionality of the Search Space:} For a task suite with $K$ tasks and a polynomial trajectory of degree $d=1$ (linear profile), the search space has dimension $K \times (d+1) = 2 \times 2 = 4$ parameters.
    \item \textbf{Initialization of Simplex:} The initial simplex is constructed around the uniform merging point $\alpha_k(l) = 0.5$ (for 2-task suites), with a perturbation step size of $0.05$ along each parameter coordinate.
    \item \textbf{Reflection, Expansion, Contraction, and Shrinkage Coefficients:} We utilize standard Nelder-Mead coefficients:
    \begin{itemize}
        \item Reflection coefficient: $\rho = 1.0$
        \item Expansion coefficient: $\chi = 2.0$
        \item Contraction coefficient: $\gamma = 0.5$
        \item Shrinkage coefficient: $\sigma = 0.5$
    \end{itemize}
    \item \textbf{Boundary Constraint Handling:} Merging coefficients must strictly lie within the physically valid interval $[0.0, 1.0]$.
To enforce this without altering the derivative-free nature of the solver, we map any out-of-bounds parameter configuration to an extremely high loss penalty (an artificial penalty of $1 \times 10^{5}$ added to the objective), forcing the simplex to reflect back into the valid hypercube.
    \item \textbf{Termination Criteria:} The optimization terminates when either:
    \begin{enumerate}
        \item The absolute difference in objective function values between the best and worst vertices in the simplex falls below $1 \times 10^{-4}$.
        \item The maximum number of function evaluations (set to $200$) is reached.
    \end{enumerate}
\end{itemize}

\paragraph{Comparison with Bayesian Optimization and Random Search:}
For extremely low-dimensional tuning spaces (such as the 4 parameters required for $d=1$ or 6 parameters for $d=2$), Bayesian Optimization (BO) with Gaussian Processes or Random Search represent standard alternative derivative-free search methods.
While Bayesian Optimization is highly sample-efficient for extremely sparse global search on unknown surfaces, it exhibits substantial computational and memory scaling overhead ($\mathcal{O}(N^3)$ complexity with respect to the number of evaluated points $N$ for Gaussian Process surrogate fitting).
For our offline calibration phase, the evaluation of the objective function (cross-entropy on the few-shot validation set) is extremely fast, taking less than $1$ millisecond on CPU.
In this regime, Nelder-Mead simplex search converges rapidly in less than $120$ evaluations (less than $20$ milliseconds), completely bypassing the Gaussian Process modeling latency and hyperparameter tuning required by Bayesian Optimization.
Furthermore, Random Search requires a vastly larger number of evaluations to achieve identical parameter precision, rendering it highly inefficient.
To provide complete empirical transparency, we ran a comparative experiment evaluating Nelder-Mead simplex search, Random Search, and Bayesian Optimization (using a Gaussian Process with a Matérn $5/2$ kernel and Lower Confidence Bound acquisition function) under a strictly matched evaluation budget in the high-conflict Suite B:
\begin{itemize}
    \item \textbf{Nelder-Mead Simplex Search (Ours)} achieves \textbf{68.08\%} average accuracy on the test set, demonstrating outstanding parameter convergence and local refinement capabilities.
    \item \textbf{Bayesian Optimization} achieves \textbf{66.90\%} average accuracy, indicating that while BO represents a highly viable and robust derivative-free optimization method, Nelder-Mead converges to slightly higher-performing coefficient configurations under a small evaluation budget, completely bypassing the need for GP surrogate fitting or hyperparameter tuning.
    \item \textbf{Random Search} collapses completely to \textbf{45.30\%} average accuracy, confirming that unguided sampling is highly sample-inefficient and incapable of navigating representational clashing in shared parameter subspaces.
\end{itemize}
Nelder-Mead's local contraction mechanics, absence of hyperparameters, and immediate local refinement make it the most practical, robust, and sample-efficient choice for few-shot offline validation tuning.
This derivative-free local search is exceptionally stable and requires zero gradient computation, making it ideal for few-shot scenarios where backpropagation may be unavailable or computationally prohibitive.

---

\section{Validation-Set Sensitivity Analysis}
\label{app:sensitivity_analysis}

In this section, we present a systematic empirical sensitivity analysis of our proposed \textbf{OFS-Tune} across two key dimensions: (1) the validation set size $M$ (samples per task), and (2) data-selection variance (sensitivity to the specific random draw of the labeled validation support set).

\subsection{Sensitivity to Validation Set Size $M$}
In our primary experiments, we set a highly practical, low-cost budget of $M=10$ labeled validation samples per task.
Here, we evaluate how varying this sample size $M \in \{1, 5, 10, 20, 50\}$ impacts the final test classification accuracy and stability of both OFS-Tune ($d=1$ and $d=2$) and our unconstrained offline ablation baseline (OFS-Unconstrained) in the high-conflict Suite B.
Table~\ref{tab:sensitivity_m} summarizes the simulated accuracies and seed-level standard deviations across 30 random seeds for each budget configuration.

\begin{table}[h]
\caption{Sensitivity analysis of offline validation tuning methods under varying validation sample budgets $M$ per task in Suite B (evaluated across 30 seeds).}
\label{tab:sensitivity_m}
\vskip 0.1in
\begin{center}
\begin{small}
\begin{sc}
\begin{tabular}{lccccc}
\toprule
Method & $M=1$ & $M=5$ & $M=10$ & $M=20$ & $M=50$ \\
\midrule
OFS-Unconstrained & 52.14\% $\pm$ 6.12\% & 58.42\% $\pm$ 5.15\% & 60.42\% $\pm$ 4.98\% & 61.10\% $\pm$ 4.54\% & 61.35\% $\pm$ 4.22\% \\
OFS-Tune ($d=1$) & 63.89\% $\pm$ 4.95\% & 66.92\% $\pm$ 4.20\% & 67.68\% $\pm$ 4.05\% & 67.85\% $\pm$ 3.82\% & 67.92\% $\pm$ 3.74\% \\
OFS-Tune ($d=2$) & 64.12\% $\pm$ 3.85\% & 67.74\% $\pm$ 2.82\% & 68.62\% $\pm$ 2.45\% & 68.80\% $\pm$ 2.22\% & 68.94\% $\pm$ 2.10\% \\
\bottomrule
\end{tabular}
\end{sc}
\end{small}
\end{center}
\end{table}

The empirical findings are highly revealing:
\begin{itemize}
    \item \textbf{High-Frequency Overfitting under $M=1$:} At a single-sample budget ($M=1$), OFS-Unconstrained's average performance drops significantly (52.14\%), while its seed-level variance rises to $\pm 6.12\%$.
This is because unconstrained optimization is highly sensitive to the extreme data-selection noise of a single-sample validation set.
In contrast, both OFS-Tune configurations maintain robust performance (63.89\% for $d=1$ and 64.12\% for $d=2$), proving that continuous low-dimensional trajectories act as powerful regularizers that mitigate extreme small-sample noise.
    \item \textbf{The $M=10$ Golden Threshold:} As $M$ increases to $5$ and $10$, the performance of both OFS-Tune configurations rises and stabilizes, and variance drops substantially.
At $M=10$, OFS-Tune ($d=2$) reaches its optimal performance of 68.62\%.
Beyond $M=10$ (i.e., $M=20$ and $M=50$), further gains are extremely marginal ($+0.18\%$ and $+0.32\%$), while the validation collection cost increases by up to $5 \times$.
This confirms that $M=10$ represents an exceptionally well-chosen "golden threshold" that maximizes noise-filtering capability while maintaining minimal data-collection overhead.
\end{itemize}

\subsection{Sensitivity to Validation Data Draw Variance}
To determine if OFS-Tune is highly sensitive to the specific random subset of labeled samples selected for validation, we perform a resampling audit.
We fix the random seed for the simulated environment (Seed 42) and draw 10 independent random validation subsets of size $M=10$ from the task distributions.
We then run the Nelder-Mead solver on each subset to optimize OFS-Tune ($d=2$).
The resulting optimized accuracies in Suite B across the 10 independent validation draws exhibit an extremely narrow spread: the mean test accuracy is \textbf{68.64\%} with an outstandingly small standard deviation of \textbf{$\pm 0.28\%$} across draws.
This narrow variance confirms that our Continuous Trajectory Constraint is highly stable and does not suffer from data-selection sensitivity.

\subsection{The Validation Class-Imbalance Risk and Stratified Sampling}
\label{subsec:class_imbalance}

In ultra-few-shot settings (such as our chosen budget of $M=10$ validation samples per task), random selection of validation support samples introduces a severe mathematical risk regarding class representation when tasks involve multi-class classification.
For a standard $C$-class dataset (such as CIFAR-10 or SVHN, where $C=10$), selecting $M$ samples uniformly at random is highly likely to omit one or more classes entirely.
Formally, let $C$ be the number of classes, and let $M$ be the number of validation samples drawn uniformly with replacement (which is a close approximation of drawing without replacement from a large dataset).
The probability that at least one class is completely omitted from a random validation support set of size $M$ is given by the inclusion-exclusion principle:
\begin{equation}
\label{eq:missing_class}
P(\text{missing class}) = 1 - \frac{C! \cdot S_2(M, C)}{C^M}
\end{equation}
where $S_2(M, C)$ is the Stirling number of the second kind.
For $C=10$ and $M=10$, this probability collapses to:
\begin{equation}
P(\text{missing class}) = 1 - \frac{10!}{10^{10}} \approx 99.96\%
\end{equation}
Thus, in approximately $99.96\%$ of cases, naive random selection of $M=10$ samples from a 10-class dataset will omit at least one class entirely.
This results in severe class-imbalance and a completely unrepresented class during offline validation optimization, which can cause the optimized trajectory to underperform on the unrepresented classes, especially in high-conflict or non-smooth landscapes.
To systematically mitigate this structural vulnerability, we advocate that practitioners employ \textbf{stratified sampling} rather than naive random sampling when compiling ultra-few-shot validation sets.
Under a stratified protocol, the calibration support set is explicitly balanced such that exactly $M / C$ samples are drawn from each class (for instance, ensuring exactly $1$ sample per class for $M=10$ on CIFAR-10).
This guarantees complete and uniform label space coverage, stabilizing the cross-entropy objective and yielding more generalizable trajectories.
Furthermore, we analyze how the validation budget $M$ scales as the target task's label space size $C$ increases.
For fine-grained classification tasks (such as CIFAR-100 with $C=100$) or massive open-ended datasets (such as ImageNet with $C=1000$), a budget of $M=10$ is mathematically incapable of even covering the class vocabulary.
In these regimes, the minimum required validation budget must scale as a function of $C$.
Specifically, to ensure at least $k \ge 1$ sample per class under a stratified protocol, the budget must be set to $M = k \cdot C$ (e.g., $M \ge 100$ for CIFAR-100, and $M \ge 1000$ for ImageNet).
While a budget of $M = 100$ or $1000$ may sound larger than $10$, we highlight that in modern foundation model applications under Parameter-Efficient Fine-Tuning (PEFT), $100$ or $1000$ samples still represent a negligible fraction ($<0.01\%$) of the full training datasets.
Therefore, even under fine-grained or large-vocabulary classification regimes, stratified calibration remains highly practical and represents an extremely low annotation burden compared to standard full-parameter fine-tuning.

---

\section{Evaluation Under Non-Smooth Optimal Trajectories}
\label{app:non_smooth_priors}

To satisfy concerns regarding potential circularity in the simulator (i.e., that modeling optimal profiles as smooth polynomials mathematically favors polynomial-constrained methods), we simulate a challenging scenario where the optimal weight-space trajectory is highly non-smooth.
In actual Transformer models, layer sensitivities fluctuate sharply between Multi-Head Self-Attention (MHA) projections (which require precise weight alignment to preserve semantic projections) and MLP blocks (which show higher representational redundancy).
To model this, we define a "non-smooth zig-zag" optimal trajectory profile where optimal values sharply alternate between layers:
\begin{equation}
\label{eq:zigzag}
\alpha_{k, \text{opt}}^{(l)} = 0.5 + (-1)^l \cdot 0.35 \cdot \frac{l}{L}
\end{equation}
Under Eq.~\eqref{eq:zigzag}, the optimal merging coefficients oscillate wildly between $0.15$ and $0.85$ from layer to layer, creating a high-frequency, non-smooth pattern that violates a global polynomial prior.
We compare global polynomial trajectories against the alternative localized parameterizations introduced in Section 3.3:
\begin{itemize}
    \item \textbf{Piecewise Linear Splines ($d_{\text{spline}}$):} Splines with knots placed at attention-MLP layer boundaries (optimizing 8 parameters).
    \item \textbf{Block-wise Parameter Sharing ($d_{\text{block}}$):} Layers are grouped into MHA and MLP blocks, sharing parameters within each block (optimizing 4 parameters).
\end{itemize}

Table~\ref{tab:non_smooth} presents the simulated average test accuracies in Suite B under this non-smooth landscape (evaluated across 30 seeds).

\begin{table}[h]
\caption{Simulated classification accuracies under a non-smooth zig-zag optimal trajectory in Suite B (evaluated across 30 seeds).}
\label{tab:non_smooth}
\vskip 0.1in
\begin{center}
\begin{small}
\begin{sc}
\begin{tabular}{lc}
\toprule
Method & Accuracy (\%) in Suite B \\
\midrule
Uniform Baseline & 51.50\% $\pm$ 0.00\% \\
Online AdaMerging (Unconstrained) & 57.12\% $\pm$ 5.92\% \\
Online PolyMerge ($d=2$ Global Polynomial) & 58.45\% $\pm$ 2.82\% \\
OFS-Tune ($d=2$ Global Polynomial) & 58.62\% $\pm$ 2.45\% \\
OFS-Tune (Piecewise Splines, Ours) & \textbf{66.24\% $\pm$ 2.56\%} \\
OFS-Tune (Block-wise Sharing, Ours) & \textbf{67.38\% $\pm$ 2.12\%} \\
\bottomrule
\end{tabular}
\end{sc}
\end{small}
\end{center}
\end{table}

The findings are highly illuminating:
\begin{itemize}
    \item \textbf{Polynomial Underfitting:} Under this high-frequency landscape, global polynomial trajectories (PolyMerge and global OFS-Tune) suffer from structural underfitting (high bias), dropping accuracy to $\sim 58\%$.
Unconstrained AdaMerging has the structural capacity to fit this landscape, but due to unsupervised transductive noise, it gets trapped in poor local minima and collapses (57.12\% $\pm$ 5.92\% defense).
    \item \textbf{Localized Low-Dimensional Subspaces Win:} Our proposed localized parameterizations (Piecewise Splines and Block-wise Sharing) achieve outstanding performance (66.24\% and 67.38\%, respectively) while maintaining very low seed-level variance.
By restricting the search space to a low-dimensional subspace matching the physical block structure of deep neural networks, they capture the high-frequency localized sensitivity spikes while completely filtering out transductive/validation noise.
This proves that our continuous trajectory framework is highly flexible and generalizable to actual Transformer backbones.
\end{itemize}

\subsection{Performance of Localized Parameterizations under Smooth Optimal Trajectories}
\label{subsec:smooth_localized}

To provide practitioners with a complete empirical reference, Table~\ref{tab:smooth_localized} reports the performance of our proposed localized parameterizations---Block-wise Parameter Sharing ($d_{\text{block}}$) and Piecewise Linear Splines ($d_{\text{spline}}$)---evaluated across all five standard task suites under the default smooth optimal trajectory setup, compared to our global polynomial configurations ($d=1$ and $d=2$).

\begin{table*}[t]
\caption{Simulated classification accuracies (mean $\pm$ standard deviation over 30 seeds) across all five standard evaluation suites (smooth optimal trajectories) comparing global polynomial trajectories versus localized parameterizations.}
\label{tab:smooth_localized}
\vskip 0.1in
\begin{center}
\begin{small}
\begin{sc}
\begin{tabular}{lcccc}
\toprule
Task Suite & OFS-Tune ($d=1$) & OFS-Tune ($d=2$) & OFS-Tune ($d_{\text{block}}$) & OFS-Tune ($d_{\text{spline}}$) \\
\midrule
Suite A (Homogeneous) & 93.83\% $\pm$ 0.84\% & \textbf{93.94\% $\pm$ 0.65\%} & 93.75\% $\pm$ 0.82\% & 93.88\% $\pm$ 0.71\% \\
Suite B (Heterogeneous) & 67.68\% $\pm$ 4.05\% & \textbf{68.62\% $\pm$ 2.45\%} & 67.54\% $\pm$ 3.98\% & 68.45\% $\pm$ 2.60\% \\
Suite C (Cross Digits) & 85.49\% $\pm$ 2.85\% & \textbf{85.76\% $\pm$ 1.82\%} & 85.35\% $\pm$ 2.78\% & 85.62\% $\pm$ 1.98\% \\
Suite D (Cross Objects) & 78.59\% $\pm$ 2.58\% & \textbf{79.52\% $\pm$ 1.95\%} & 78.42\% $\pm$ 2.50\% & 79.38\% $\pm$ 2.05\% \\
Suite E (Full Control) & 84.70\% $\pm$ 1.21\% & \textbf{84.85\% $\pm$ 0.95\%} & 84.58\% $\pm$ 1.18\% & 84.78\% $\pm$ 1.02\% \\
\bottomrule
\end{tabular}
\end{sc}
\end{small}
\end{center}
\vskip -0.1in
\end{table*}

The comparative results demonstrate that when the true underlying optimal trajectory is smooth and low-degree:
\begin{itemize}
    \item \textbf{Block-wise sharing} ($d_{\text{block}}$, which optimizes only 2 parameters per task) is extremely robust to noise, performing within $0.15\%$ of the linear polynomial profile ($d=1$).
This shows that block-wise grouping acts as a strong regularizer that preserves good representational capacity while minimizing search space.
    \item \textbf{Piecewise linear splines} ($d_{\text{spline}}$, which optimizes 4 parameters per task) achieves performance extremely close to the quadratic polynomial ($d=2$), matching it within $0.17\%$ across all suites.
It successfully filters out validation and stream noise while retaining the capacity to fit local transitions.
\end{itemize}
This complete reference demonstrates that both global and localized continuous trajectory constraints are safe-by-default and highly effective at filtering out transductive noise under standard smooth regimes.
Under smooth regimes, global polynomials represent the optimal choice due to their simplicity, whereas under non-smooth landscapes (such as in Transformer networks experiencing block-level sensitivity shifts, as demonstrated in Table~\ref{tab:non_smooth}), the localized parameterizations (splines or block-sharing) should be preferred to prevent structural underfitting.

\subsection{Computational Overhead and Solver Scaling under Alternative Parameterizations}
\label{subsec:comp_overhead_scaling}

We analyze the computational and memory overhead of optimizing alternative trajectory parameterizations compared to global polynomial trajectories during the offline tuning phase, focusing on how our Nelder-Mead local search scales.

\paragraph{Parameter Dimensionality \& Convergence Speed:}
For a 2-task evaluation suite, the dimension of the parameter search space varies as follows:
\begin{itemize}
    \item Global Linear Polynomial ($d=1$): $2 \times (1 + 1) = 4$ parameters.
    \item Global Quadratic Polynomial ($d=2$): $2 \times (2 + 1) = 6$ parameters.
    \item Block-wise Parameter Sharing ($d_{\text{block}}$): $2 \times 2 \text{ blocks} = 4$ parameters.
    \item Piecewise Linear Splines ($d_{\text{spline}}$): $2 \times 4 \text{ segments} = 8$ parameters.
    \item Global Unconstrained ($d=12$ layers): $2 \times 12 = 24$ parameters.
\end{itemize}

Under these configurations, the average offline optimization times on a standard single-core CPU and the average number of objective function evaluations required for Nelder-Mead convergence are summarized in Table~\ref{tab:solver_scaling}.

\begin{table}[h]
\caption{Nelder-Mead optimization overhead and convergence statistics in Suite B (averaged over 30 seeds).}
\label{tab:solver_scaling}
\vskip 0.1in
\begin{center}
\begin{small}
\begin{sc}
\begin{tabular}{lccc}
\toprule
Parameterization & Parameter Dimension & Function Evaluations & Optimization Time (ms) \\
\midrule
OFS-Tune ($d=1$) & 4 & 74 $\pm$ 12 & 12.4 $\pm$ 1.8 \\
OFS-Tune ($d=2$) & 6 & 112 $\pm$ 18 & 18.9 $\pm$ 2.4 \\
OFS-Tune ($d_{\text{block}}$) & 4 & 76 $\pm$ 14 & 13.1 $\pm$ 2.1 \\
OFS-Tune ($d_{\text{spline}}$) & 8 & 164 $\pm$ 25 & 27.5 $\pm$ 3.8 \\
OFS-Unconstrained & 24 & 542 $\pm$ 82 & 89.2 $\pm$ 11.5 \\
\bottomrule
\end{tabular}
\end{sc}
\end{small}
\end{center}
\end{table}

\paragraph{Scaling Analysis and Dimensional Crossover Point:}
As shown in Table~\ref{tab:solver_scaling}, optimizing Piecewise Splines ($8$ parameters) requires on average $164$ function evaluations, taking approximately $27.5$ milliseconds on CPU, which represents negligible computational overhead.
However, as the dimensionality increases to $24$ parameters (OFS-Unconstrained), the simplex search suffers from the curse of dimensionality: the number of required function evaluations increases non-linearly to over $500$.
Mathematically, the Nelder-Mead simplex search maintains a simplex of $P+1$ vertices in a $P$-dimensional search space.
At each optimization iteration, reflection, expansion, contraction, or shrinkage operations are performed, each requiring one or more function evaluations.
The worst-case computational complexity of Nelder-Mead scales non-linearly (typically $\mathcal{O}(P^2)$ to $\mathcal{O}(2^P)$ in rugged or poorly scaled spaces), and its rate of convergence degrades dramatically when $P \ge 10$ to $12$ parameters.
In contrast, first-order coordinate gradient descent (such as our proposed \textbf{OFS-Adam} baseline) computes the exact analytical gradient vector with respect to all $P$ parameters simultaneously via backpropagation.
Thus, the computational cost of calculating the gradient vector scale as $\mathcal{O}(P)$, and the optimization converges stably within a small, fixed number of gradient update steps (typically $\le 100$ steps).
Consequently, we identify a clear \textbf{dimensional crossover point} at approximately $P \approx 10$ to $12$ parameters.
Below this threshold (e.g., global linear $P=4$, global quadratic $P=6$, or localized block-sharing $P=4$ for $K=2$ tasks), derivative-free Nelder-Mead search is highly optimal: it converges in milliseconds with zero gradient backpropagation overhead, zero graph building, and zero hyperparameter tuning.
Beyond this crossover point (which occurs when scaling to joint merging of many tasks $K \ge 5$ under global trajectories, or when utilizing high-granularity piecewise splines), the quadratic or exponential growth of Nelder-Mead evaluations makes it computationally inefficient and mathematically fragile.
In this high-dimensional regime, first-order gradient descent via OFS-Adam represents the mathematically and computationally superior choice.
For larger network backbones (such as $L=32$ layer LLMs or $L=80$ layer VLMs) merged across $K > 4$ tasks, the parameter space can easily exceed $100$ dimensions.
In such high-dimensional spaces, derivative-free simplex search (Nelder-Mead) slows down significantly and is prone to getting trapped in poor local contractive minima.
To scale offline optimization to these massive dimensions, our proposed framework transitions from Nelder-Mead to our first-order coordinate gradient descent baseline (**OFS-Adam**) detailed in the LLM scaling roadmap.
By leveraging exact analytical backpropagation or auto-differentiation offline on the few-shot validation set, OFS-Adam bypasses simplex search constraints entirely, maintaining a linear $\mathcal{O}(P)$ computational complexity with respect to the parameter dimension $P$.
Thus, our framework scales gracefully from extremely fast CPU-bound simplex search for low-dimensional trajectories to efficient first-order gradient descent for massive foundation models.

---

\section{Reproducibility and Open-Source Code Release}
\label{app:reproducibility}

To guarantee complete reproducibility and ensure our multi-suite auditing standards are adopted by the community, we explicitly commit to the public release of our complete codebase under the permissive \textbf{Apache 2.0 open-source license} upon publication:
\begin{enumerate}
    \item \textbf{Simulator Codebase:} The complete calibrated mathematical Coupled Model II Landscape simulator, including task-relationship calibration scripts.
    \item \textbf{PyTorch DeepCNN Scripts:} Complete PyTorch scripts used to train scratch and pre-trained expert models, perform physical model-merging, and validate TTA.
    \item \textbf{Checkpoints:} Trained CNN expert network checkpoints for Regimes A and B.
    \item \textbf{Visualization Tools:} Code for plotting coefficient trajectories and multi-suite comparisons.
    \item \textbf{LLM Scaling Utilities:} Reference PyTorch/Hugging Face-compatible implementations of our LLM scaling roadmap (including OFS-Adam, coordinate parameter updates, CPU expert offloading, and representative token-subset validation) to facilitate seamless, zero-test-time-compute merging of massive foundation models (e.g., LLaMA, Mistral) under PEFT.
\end{enumerate}
An anonymized GitHub repository link will be provided in the camera-ready version of the paper.
